\chapter*{General Introduction}
\addcontentsline{toc}{chapter}{General Introduction}

The digital transformation of political communication has produced an unprecedented volume of public information. News agencies, parliamentary bodies, governmental institutions, research organizations, and civil society actors publish large numbers of articles, reports, and official documents every day. These sources contain valuable signals about actors, events, institutions, alliances, conflicts, financial links, and temporal changes, but most of this information remains embedded in unstructured text.

For analysts, journalists, and researchers, the difficulty is no longer only to access information, but to connect dispersed facts and interpret them in a reliable way. Keyword search can retrieve documents, and classical retrieval-augmented generation can summarize passages, yet both approaches remain limited when the question requires entity disambiguation, multi-hop reasoning, temporal filtering, and explicit source attribution.

This difficulty is amplified by the rhythm of political information. Relevant evidence may appear first in an RSS entry, then in a full news article, then in an institutional report or policy PDF. A platform designed for political intelligence must therefore do more than consume a static dataset. It must continuously discover, collect, normalize, and evaluate heterogeneous sources while preserving the provenance of every fact. In this project, that role is fulfilled by an autonomous crawler composed of several complementary levels: curated RSS monitoring, query-driven news discovery, PDF report ingestion, and administrative controls for scheduled collection.

Within its research and development activities, SCIENTIAE is developing a
sovereign Graph-RAG framework designed to adapt to different domains rather
than remain tied to a single use case. The framework provides reusable
components for data acquisition, LLM-based knowledge extraction, graph
construction, semantic retrieval, evidence-grounded answer generation, and
system evaluation. A new domain can therefore be integrated by adapting its
sources, entity and relation models, extraction rules, and evaluation
questions while preserving the same core architecture.

Politics and geopolitics were selected as the application and validation domain for this work because they provide a demanding environment for testing the framework. Political intelligence analysis often depends on relational questions: a user may need to understand how two actors are connected, which events created a diplomatic shift, how an institution evolved over time, or which documents support a specific claim. Such questions cannot be answered robustly by treating documents as isolated text fragments; they require a representation that combines textual evidence with structured relations.

The implemented system constitutes a domain-specific instance of this broader
multimodal Graph-RAG framework. It combines dense semantic retrieval, a Neo4j
knowledge graph, DSPy-optimized entity and relation extraction during
ingestion, and an LLM-based reasoning layer. Its goal is to transform a large
corpus of political and geopolitical information into an exploitable
knowledge layer that supports grounded, explainable, and temporally aware
analysis. The resulting platform serves both as a functional decision-support
tool and as an experimental environment for assessing how the framework can
later be adapted to other domains.

A second challenge concerns data acquisition itself. Many Graph-RAG systems assume that the corpus already exists and focus only on retrieval and generation. In this project, the corpus is produced by the platform. The crawler must avoid duplicates, preserve source metadata, handle heterogeneous formats, operate politely with external websites, and provide enough operational visibility for an administrator to supervise ingestion. This makes crawling and ingestion a central contribution of the project rather than a peripheral utility.

The final challenge is trust. Political analysis is sensitive because a generated answer may influence the interpretation of institutions, actors, or conflicts. The system must therefore expose its evidence, separate raw documents from extracted relations, preserve timestamps, and make retrieval behavior auditable. This need explains the combination of MongoDB for document storage, Neo4j for explicit relations, ZVec for local semantic search, and observability tools for operational monitoring.

To address these challenges, the project pursues three main objectives. The
first is to contribute to a modular Graph-RAG framework whose principal
processing stages can be configured for different domains. The second is to
instantiate this framework for political and geopolitical analysis by
implementing an autonomous ingestion and knowledge-graph construction
pipeline together with a reasoning engine that combines vector retrieval,
graph traversal, and cited answer synthesis. The third is to evaluate the
system in terms of extraction reliability, retrieval quality, answer
grounding, reasoning depth, and response time.

These objectives are translated into concrete engineering goals. The platform must collect data autonomously from RSS feeds, Google News queries, direct URLs, and PDF sources. It must normalize the collected material into document records that can be processed by the same extraction pipeline. It must resolve entities deterministically through canonical names and stable identifiers, without depending on a Wikidata backbone as the main source of truth. It must also provide an administration interface where ingestion status, scheduler state, logs, source filters, and backfill operations can be controlled without modifying code.

From a user perspective, the objective is to provide an integrated analytical workspace. The analyst can ask questions in natural language, inspect sources, explore graph neighborhoods, follow political trends, consult documents, and observe system health. From an engineering perspective, the objective is to build a modular pipeline in which crawling, extraction, graph publication, chunking, embedding, retrieval, and evaluation remain separable and testable.

The expected contribution is therefore both methodological and practical. Methodologically, the project studies how autonomous crawling, knowledge graph construction, and retrieval-augmented generation can be combined in a reusable, domain-adaptable workflow. Practically, it delivers and evaluates a political and geopolitical implementation in which source acquisition, evidence structuring, semantic search, graph exploration, and answer generation are connected through operational backend services and frontend controls.

The remainder of this report is organized into five chapters. Chapter~1
introduces the host organization, project context, problem statement,
objectives, actors, requirements, technology choices, and project-management
approach. Chapter~2 presents the state of the art concerning knowledge graphs,
retrieval-augmented generation, Graph-RAG, ingestion strategies, embeddings,
and large language models. Chapter~3 describes the system architecture, data
flows, domain model, ingestion pipeline, and reasoning design. Chapter~4
details the implementation, persistence layers, interfaces, deployment,
observability, and optimization mechanisms. Finally, Chapter~5 presents the
system evaluation and performance analysis, including RAGAS metrics,
reasoning-mode comparison, embedding performance, DSPy-based knowledge
extraction optimization, requirements validation, and limitations.
